\section{Measures of Dispersion}

\subsection{Dispersion or Variation}
While measures of central tendency tell us around what values a dataset tends to cluster, \emph{measures of dispersion} give us an idea of how spread out the data points are from one another and from the mean.

\begin{remark}
	Two datasets can have the exact same mean but very different dispersion. For example, the sets $\{49, 50, 51\}$ and $\{0, 50, 100\}$ have the same mean ($\bar{X} = 50$), but clearly the second set is much more dispersed.
\end{remark}

Below we examine some of the most commonly used measures of dispersion in statistics.

\subsection{Range}
The \emph{range} of a dataset is the difference between the largest and smallest values in the set.

\begin{example}
	The range of the set $2,3,3,5,5,5,8,10,12$ is $12-2=10.$
\end{example}

\begin{remark}
	The range is the simplest measure of dispersion, but it has the disadvantage of depending solely on the extreme values, ignoring the distribution of intermediate data points.
\end{remark}


\subsection{Mean Deviation}
The \emph{mean deviation} or \emph{average deviation} of a set of $N$ numbers $X_{1},...,X_{N}$ is defined as
\begin{align}
	DM=\dfrac{\sum \abs{X_{j}-\bar{X}}}{N}
\end{align}
where $\bar{X}$ is the arithmetic mean of the numbers and $\abs{\cdot}$ denotes absolute value.

\begin{example}
	Find the mean deviation of the list $2,3,6,8,11.$
\end{example}

\begin{solution}
	First, we calculate the mean:
	\begin{align}
		\bar{X} = \frac{2+3+6+8+11}{5} = \frac{30}{5} = 6
	\end{align}
	
	Then we calculate the absolute deviations from the mean:
	\begin{align}
		DM &= \frac{|2-6| + |3-6| + |6-6| + |8-6| + |11-6|}{5} \\
		&= \frac{4 + 3 + 0 + 2 + 5}{5} = \frac{14}{5} = 2.8
	\end{align}
\end{solution}


\subsection{Standard Deviation}
The \emph{standard deviation} of a set of $N$ numbers $X_{1},...,X_{N}$ is denoted by $s$ and is defined as
\begin{align}
	s=\sqrt{\dfrac{\sum\left( X_{j}-\bar{X} \right)^{2}}{N}}=\sqrt{\sum\dfrac{x_{j}^{2}}{N}}
\end{align} where $x_{j}:=X_{j}-\bar{X}.$

\begin{remark}
	The standard deviation measures how spread out the data are from the mean. A small standard deviation indicates that data points cluster tightly around the mean, whereas a large standard deviation indicates that data points are widely spread out.
\end{remark}

If $X_{1},...,X_{N}$ occur with frequencies $f_{1},...,f_{N}$ respectively, the standard deviation can be expressed as
\begin{align}
	s=\sqrt{\dfrac{\sum f_{j}\left( X_{j}-\bar{X} \right)^{2}}{N}}=\sqrt{\dfrac{\sum f_{j}x_{j}^{2}}{N}}
\end{align}

\begin{example}
	Calculate the standard deviation of the set $2, 4, 4, 4, 5, 5, 7, 9.$
\end{example}

\begin{solution}
	First, we calculate the mean:
	\begin{align}
		\bar{X} = \frac{2+4+4+4+5+5+7+9}{8} = \frac{40}{8} = 5
	\end{align}
	
	Then we calculate the squared deviations:
	\begin{align}
		(2-5)^2 &= 9 \\
		(4-5)^2 &= 1 \quad (\text{three times}) \\
		(5-5)^2 &= 0 \quad (\text{two times}) \\
		(7-5)^2 &= 4 \\
		(9-5)^2 &= 16
	\end{align}
	
	Therefore:
	\begin{align}
		s &= \sqrt{\frac{9 + 3(1) + 2(0) + 4 + 16}{8}} = \sqrt{\frac{32}{8}} = \sqrt{4} = 2
	\end{align}
\end{solution}


\subsection{Variance}
The \emph{variance} of a set of numbers is defined as the square $s^{2}$ of the standard deviation $s$.

\begin{remark}
	In statistics, it is important to distinguish between the standard deviation of a \emph{population} and that of a \emph{sample}. To distinguish them, we use $\sigma$ for the population and continue using $s$ for the sample.
	
	When working with a sample and estimating the population variance, we use $N-1$ instead of $N$ in the denominator:
	\begin{align}
		s^2 = \frac{\sum (X_j - \bar{X})^2}{N-1}
	\end{align}
	
	This adjustment (called \emph{Bessel's correction}) yields an unbiased estimator of the population variance. For large samples ($N > 30$), the difference between using $N$ and $N-1$ is virtually negligible.
\end{remark}


\subsection{Shortcut Method}
A computationally more efficient way to calculate the variance is:
\begin{align}
	s^{2} &= \overline{X^{2}} - \overline{X}^{2} \\
	&= \frac{\sum X_j^2}{N} - \left(\frac{\sum X_j}{N}\right)^2
\end{align}

where $\overline{X^2}$ denotes the mean of the squares and $\overline{X}^2$ denotes the square of the mean.

\begin{remark}
	This formula is equivalent to the original definition but avoids calculating individual deviations. It is particularly useful for manual calculations or when working with grouped data.
\end{remark}


\subsection{Coefficient of Variation}
The \emph{coefficient of variation} is a relative measure of dispersion that allows us to compare the variability of datasets measured in different units or scales:
\begin{align}
	CV = \frac{s}{\bar{X}} \times 100\%
\end{align}

\begin{example}
	Compare the variability of heights and weights in a group of individuals:
	\begin{itemize}
		\item Heights: $\bar{X} = 170$ cm, $s = 8$ cm
		\item Weights: $\bar{X} = 70$ kg, $s = 10$ kg
	\end{itemize}
\end{example}

\begin{solution}
	\begin{align}
		CV_{\text{height}} &= \frac{8}{170} \times 100\% \approx 4.71\% \\
		CV_{\text{weight}} &= \frac{10}{70} \times 100\% \approx 14.29\%
	\end{align}
	
	Even though the standard deviation of weight ($10$ kg) is greater than that of height ($8$ cm) in absolute terms, the coefficient of variation reveals that weight exhibits much greater relative variability ($14.29\%$ vs $4.71\%$).
\end{solution}


\subsection{Variance of Combined Sets}
If two sets of $N_{1}$ and $N_{2}$ observations have variances $s_{1}^{2}$ and $s_{2}^{2}$ respectively but share the same arithmetic mean $\bar{X},$ then the variance of the combined set is
\begin{align}
	s^{2}=\dfrac{N_{1}s_{1}^{2}+N_{2}s_{2}^{2}}{N_{1}+N_{2}}.
\end{align}


\subsection{Chebyshev's Theorem}

\begin{theorem}[Chebyshev]
	For any random variable $X$ with mean $\mu$ and standard deviation $\sigma$, and for any real number $k > 1$,
	\begin{align}
		P(|X - \mu| < k\sigma) \geq 1 - \frac{1}{k^2}
	\end{align}
	
	That is, at least $1 - \frac{1}{k^2}$ of the probability distribution of $X$ lies within $k$ standard deviations of the mean.
\end{theorem}

\begin{proof}
	We will prove the theorem for the discrete case. Let $f(x)$ be the probability mass function of $X$.
	
	The variance is:
	\begin{align}
		\sigma^2 = \sum_{x} (x - \mu)^2 f(x)
	\end{align}
	
	We split the summation into two parts: values inside $k\sigma$ from the mean and values outside:
	\begin{align}
		\sigma^2 &= \sum_{|x-\mu| < k\sigma} (x-\mu)^2 f(x) + \sum_{|x-\mu| \geq k\sigma} (x-\mu)^2 f(x)
	\end{align}
	
	The first summation is non-negative, hence:
	\begin{align}
		\sigma^2 \geq \sum_{|x-\mu| \geq k\sigma} (x-\mu)^2 f(x)
	\end{align}
	
	For terms in the second summation, $(x-\mu)^2 \geq k^2\sigma^2$, so:
	\begin{align}
		\sigma^2 &\geq \sum_{|x-\mu| \geq k\sigma} k^2\sigma^2 f(x) = k^2\sigma^2 \sum_{|x-\mu| \geq k\sigma} f(x) \\
		&= k^2\sigma^2 P(|X-\mu| \geq k\sigma)
	\end{align}
	
	Dividing by $k^2\sigma^2$:
	\begin{align}
		P(|X-\mu| \geq k\sigma) \leq \frac{1}{k^2}
	\end{align}
	
	Therefore:
	\begin{align}
		P(|X-\mu| < k\sigma) = 1 - P(|X-\mu| \geq k\sigma) \geq 1 - \frac{1}{k^2}
	\end{align}
\end{proof}

\begin{remark}
	For $k=2$, the theorem guarantees that at least $1 - \frac{1}{4} = 75\%$ of the data lie within 2 standard deviations of the mean. For $k=3$, at least $1 - \frac{1}{9} \approx 88.9\%$ lie within 3 standard deviations.
	
	In normal distributions, these percentages are significantly higher: $95.45\%$ within $\pm 2\sigma$ and $99.73\%$ within $\pm 3\sigma$.
\end{remark}


\subsection{Python}

\texttt{Numpy} provides functions to calculate all measures of dispersion:

\begin{lstlisting}[language=Python]
import numpy as np

data = np.array([2, 4, 4, 4, 5, 5, 7, 9])

# Range
rango = np.ptp(data)
print(f"Range: {rango}")  # 7

# Mean deviation
desv_media = np.mean(np.abs(data - np.mean(data)))
print(f"Mean deviation: {desv_media}")  # 1.5

# Population standard deviation (ddof=0)
std_pop = np.std(data, ddof=0)
print(f"Population standard deviation: {std_pop}")  # 2.0

# Sample standard deviation (ddof=1)
std_muestra = np.std(data, ddof=1)
print(f"Sample standard deviation: {std_muestra}")  # 2.138

# Variance
var_pop = np.var(data, ddof=0)
print(f"Population variance: {var_pop}")  # 4.0

var_muestra = np.var(data, ddof=1)
print(f"Sample variance: {var_muestra}")  # 4.571

# Coefficient of variation
cv = (std_pop / np.mean(data)) * 100
print(f"Coefficient of variation: {cv:.2f}%")  # 40.00%
\end{lstlisting}

\begin{remark}
	The \texttt{ddof} (Delta Degrees of Freedom) parameter controls the denominator: \texttt{ddof=0} divides by $N$ (population) while \texttt{ddof=1} divides by $N-1$ (sample).
\end{remark}
